% GENERATED FILE. Edit canonical lessons, then run npm run generate:books.
% canonical-source-hash: fnv1a64:15c65e59dcb76257
% canonical-lessons: TE-C249-bharincu, TE-C249-paku, TE-C249-vangu, TE-C249-cacu, TE-C249-percu

\chapter{To endure, To crawl, To bend down, To stretch out, To stack}
\label{ch:te-to-endure-to-crawl-to-bend-down-to-stretch-out-to-stack}
\hlchaptermodality{249}

\begin{chapteropening}
\textbf{By the end of this chapter:} \emph{I can say to endure, to crawl, to bend down, to stretch out and to stack.}

Complete the last lesson of chapter 249: I can say to endure, to crawl, to bend down, to stretch out and to stack.
\end{chapteropening}

\section[\texorpdfstring{\te{భరించు} (bhariñcu)}{భరించు (bhariñcu)}]{\te{భరించు} (bhariñcu) — to endure, to put up with}

\label{lesson:TE-C249-bharincu}



\begin{quote}
Before the new one: say the Telugu for to rob, then the Telugu for to hand over.
\end{quote}

\subsection*{You'll want to know: \te{భరించు}}
\textbf{\te{భరించు}} — \emph{bhariñcu} — \textquotedblleft{}to endure, to put up with\textquotedblright{}.

\begin{grammarlens}[title={What you've built}]
One more word for this chapter.
\end{grammarlens}

\subsection*{Your turn}
\begin{itemize}
\raggedright
  \item \emph{Say it:} \emph{bhariñcu}
  \item \emph{Say it:} \emph{bhariñcu}, once more
  \item \emph{Say it:} say \emph{bhariñcu}
  \item \emph{Recall:} say the Telugu for to rob, then the Telugu for to hand over, then say \emph{bhariñcu} again
\end{itemize}

\begin{tcolorbox}[breakable,colback=teal!4,colframe=teal!35!black,title={Before you move on}]
What does \te{భరించు} mean? (\textquotedblleft{}To endure, to put up with\textquotedblright{}.) Say it once more.
\end{tcolorbox}
\section[\texorpdfstring{\te{పాకు} (pāku)}{పాకు (pāku)}]{\te{పాకు} (pāku) — to crawl}

\label{lesson:TE-C249-paku}



\begin{quote}
Before the new one: say the Telugu for to hand over, then the Telugu for to endure.
\end{quote}

\subsection*{You'll want to know: \te{పాకు}}
\textbf{\te{పాకు}} — \emph{pāku} — \textquotedblleft{}to crawl\textquotedblright{}.

\begin{grammarlens}[title={What you've built}]
One more word for this chapter.
\end{grammarlens}

\subsection*{Your turn}
\begin{itemize}
\raggedright
  \item \emph{Say it:} \emph{pāku}
  \item \emph{Say it:} \emph{pāku}, once more
  \item \emph{Say it:} say \emph{pāku}
  \item \emph{Recall:} say the Telugu for to hand over, then the Telugu for to endure, then say \emph{pāku} again
\end{itemize}

\begin{tcolorbox}[breakable,colback=teal!4,colframe=teal!35!black,title={Before you move on}]
What does \te{పాకు} mean? (\textquotedblleft{}To crawl\textquotedblright{}.) Say it once more.
\end{tcolorbox}
\section[\texorpdfstring{\te{వంగు} (vaṅgu)}{వంగు (vaṅgu)}]{\te{వంగు} (vaṅgu) — to bend down, to stoop}

\label{lesson:TE-C249-vangu}



\begin{quote}
Before the new one: say the Telugu for to endure, then the Telugu for to crawl.
\end{quote}

\subsection*{You'll want to know: \te{వంగు}}
\textbf{\te{వంగు}} — \emph{vaṅgu} — \textquotedblleft{}to bend down, to stoop\textquotedblright{}.

\begin{grammarlens}[title={What you've built}]
One more word for this chapter.
\end{grammarlens}

\subsection*{Your turn}
\begin{itemize}
\raggedright
  \item \emph{Say it:} \emph{vaṅgu}
  \item \emph{Say it:} \emph{vaṅgu}, once more
  \item \emph{Say it:} say \emph{vaṅgu}
  \item \emph{Recall:} say the Telugu for to endure, then the Telugu for to crawl, then say \emph{vaṅgu} again
\end{itemize}

\begin{tcolorbox}[breakable,colback=teal!4,colframe=teal!35!black,title={Before you move on}]
What does \te{వంగు} mean? (\textquotedblleft{}To bend down, to stoop\textquotedblright{}.) Say it once more.
\end{tcolorbox}
\section[\texorpdfstring{\te{చాచు} (cācu)}{చాచు (cācu)}]{\te{చాచు} (cācu) — to stretch out (a limb)}

\label{lesson:TE-C249-cacu}



\begin{quote}
Before the new one: say the Telugu for to crawl, then the Telugu for to bend down.
\end{quote}

\subsection*{You'll want to know: \te{చాచు}}
\textbf{\te{చాచు}} — \emph{cācu} — \textquotedblleft{}to stretch out (a limb)\textquotedblright{}.

\begin{grammarlens}[title={What you've built}]
One more word for this chapter.
\end{grammarlens}

\subsection*{Your turn}
\begin{itemize}
\raggedright
  \item \emph{Say it:} \emph{cācu}
  \item \emph{Say it:} \emph{cācu}, once more
  \item \emph{Say it:} say \emph{cācu}
  \item \emph{Recall:} say the Telugu for to crawl, then the Telugu for to bend down, then say \emph{cācu} again
\end{itemize}

\begin{tcolorbox}[breakable,colback=teal!4,colframe=teal!35!black,title={Before you move on}]
What does \te{చాచు} mean? (\textquotedblleft{}To stretch out (a limb)\textquotedblright{}.) Say it once more.
\end{tcolorbox}
\section[\texorpdfstring{\te{పేర్చు} (pērcu)}{పేర్చు (pērcu)}]{\te{పేర్చు} (pērcu) — to stack, to pile up}

\label{lesson:TE-C249-percu}



\begin{quote}
Before the new one: say the Telugu for to bend down, then the Telugu for to stretch out.
\end{quote}

\subsection*{You'll want to know: \te{పేర్చు}}
\textbf{\te{పేర్చు}} — \emph{pērcu} — \textquotedblleft{}to stack, to pile up\textquotedblright{}.

\begin{grammarlens}[title={What you've built}]
One more word for this chapter.
\end{grammarlens}

\subsection*{Your turn}
\begin{itemize}
\raggedright
  \item \emph{Say it:} \emph{pērcu}
  \item \emph{Say it:} \emph{pērcu}, once more
  \item \emph{Say it:} say \emph{pērcu}
  \item \emph{Recall:} say the Telugu for to bend down, then the Telugu for to stretch out, then say \emph{pērcu} again
\end{itemize}

\begin{tcolorbox}[breakable,colback=teal!4,colframe=teal!35!black,title={Before you move on}]
What does \te{పేర్చు} mean? (\textquotedblleft{}To stack, to pile up\textquotedblright{}.) Say it once more.
\end{tcolorbox}
